\documentclass[10pt]{article}
\usepackage[a4paper,margin=26mm]{geometry}
\usepackage{amsmath,amssymb,amsthm,hyperref}
\newtheorem{theorem}{Theorem}
\newtheorem{lemma}{Lemma}
\newcommand{\ddc}{dd^c}
\newcommand{\MA}{\operatorname{MA}}
\newcommand{\PP}{\mathbf P}
\title{A finite-entropy initial singularity can erase an absolute zero-mass endpoint}
\author{An unreviewed supporting candidate for Problem P19.2}
\date{6 October 2026}
\begin{document}
\maketitle
\begin{abstract}
On every complex projective space of dimension at least two, we construct
an unbounded full-mass finite-energy initial potential with finite
probability entropy. For an actual bounded nef non-Archimedean model datum,
the zero-start maximal ray has a nontrivial absolute $P$-model endpoint
of zero nonpluripolar mass, whereas the maximal ray from this singular
initial potential has identically minus-infinite zero level. Both rays
have entropy slope one and remain at one fixed finite $d_1$ distance.
The example rules out absolute endpoint invariance under arbitrary
finite-entropy initial-point changes. It does not rule out the entire
residual model approximation problem or assert a failure of relative
envelopes adapted to the prescribed initial potential.
\end{abstract}

\paragraph{Conventions and sole theorem premise.}
Let $X=\PP^n_{\mathbf C}$, $n\geq2$, $\omega=\omega_{\rm FS}$,
$\int_X\omega^n=1$, $\ddc=(i/2\pi)\partial\bar\partial$ and
$t=-\log|\tau|^2$. Set
\[
 E(u)=\frac1{n+1}\sum_{k=0}^n\int_Xu\,
        \langle\omega_u^k\rangle\wedge\omega^{n-k},\qquad
 H(u)=\operatorname{Ent}(\langle\omega_u^n\rangle\mid\omega^n).
\]
For full-mass $u$ these are probability conventions. The sole substantial
theorem premise is the complete established P19.1 mathematical theorem
in its fixed smooth compact K\"ahler signature, denoted $R$: for every
actual $w\in\mathcal E^1$ and $\phi\in\mathcal E^{1,\rm NA}$ the
actual maximal weak ray $r(w,\phi)$ exists, is ordered in both arguments,
has $E(r(w,\phi)_t)=E(w)+tE^{\rm NA}(\phi)$, and an actual weak ray
is maximal for its usual datum exactly when its energy defect vanishes.
A maximal ray with a finite-entropy slice has entropy slope
$H^{\rm NA}(\phi)$. The complete theorem body, not a historical verdict,
is the premise. It contains no prescribed endpoint invariance grant.

Write
\[
 \chi_u(\lambda)=\inf_{t\geq0}(u_t-t\lambda),\qquad
 P_\omega[v]=\left(\lim_{C\to\infty}P_\omega(v+C,0)\right)^*.
\]
An identically $-\infty$ level is assigned raw mass zero by convention;
it is not a finite-energy function and is assigned no singularity metric.
All $P$ ceilings here are absolute ceilings for the smooth background
$\omega$. No singular initial point is substituted for that background.

\begin{theorem}\label{main}
For every $n\geq2$ there exist a nonpositive
$w\in\mathcal E^1(X,\omega)$ with $H(w)<\infty$, singular only at
$p=[1:0:\cdots:0]$, and an actual bounded nef model datum $\phi$, with
\[
 E^{\rm NA}(\phi)=-\frac n{n+1},\qquad H^{\rm NA}(\phi)=1,
\]
such that $a=r(0,\phi)$ and $b=r(w,\phi)$ satisfy
\[
 \chi_a(0)=w_H=\log\frac{|Z_0|^2}{\sum_{i=0}^n|Z_i|^2},
 \quad P_\omega[w_H]=w_H,\quad
 \int_X\langle\omega_{w_H}^n\rangle=0,\qquad
 \chi_b(0)\equiv-\infty.                              \tag{1}
\]
For every finite time $d_1(a_t,b_t)=-E(w)$, and both entropy slopes
are one. In fact every finite-time entropy of both rays is finite.
\end{theorem}

\section{A full-mass isolated singularity of finite entropy}
In the affine chart $Z_0\ne0$, put $z_i=Z_i/Z_0$,
$\rho=\log\sum_{i=1}^n|z_i|^2$, $Q_0(\rho)=\log(1+e^\rho)$,
$\delta=1/3$, $\rho_B=-\log2$, and
$F(x)=x\log x+(1-x)\log(1-x)$. Set
\[
 C=1-\delta-F(\delta),\qquad
 Q(\rho)=
 \begin{cases}
 -(-\rho)^{1/3}+C,&\rho\leq-1,\\
 \delta\rho-F(\delta),&-1\leq\rho\leq\rho_B,\\
 Q_0(\rho),&\rho\geq\rho_B,
 \end{cases}
 \qquad w=Q-Q_0.                                      \tag{2}
\]
Both values and first derivatives match. The function $Q$ is convex
and increasing; thus its radial weight is psh in $\mathbf C^n$.
Outside a ball $w=0$, so (2) glues to a global $\omega$-psh function.
The middle line is a tangent to $Q_0$. For $L=-\rho\geq1$,
$Q'=\tfrac13L^{-2/3}>Q_0'(-L)$, because
$L^{2/3}e^{-L}/(1+e^{-L})$ decreases there and is less than $1/3$
at one. The difference $Q-Q_0$ increases to its nonpositive value
at $-1$. Consequently $w\leq0$ and $\sup w=0$.
Its sole pole is $p$, and $Q'(-\infty)=0$.

Here are explicit mass, energy and entropy calculations, rather than
an inference from zero Lelong number. For a radial total weight $Q$,
the radial distribution of its top Monge--Amp\`ere measure is
$d(Q'^n)$, with normalized angular measure. Its mixed distribution
with $k$ copies of $\ddc Q$ and $n-k$ of $\ddc Q_0$ is
\[
                    d\bigl(Q'^k Q_0'^{\,n-k}\bigr).     \tag{3}
\]
These formulas follow by expanding the radial and tangential Hessian
eigenvalues and polarizing. They also hold for the present $C^1$
piecewise smooth convex weight by smooth convex approximation.
Continuity of $Q'$ gives no interface mass. Since $Q'$ increases from
zero to one, the measure on $X\setminus\{p\}$ has total mass one.
Thus the nonpluripolar product has full mass.

On the inner region its exact radial top measure is
\[
 d\MA(w)=n\,3^{-n}\frac23 L^{-2n/3-1}\,dL,
 \qquad |w|=O(1+L^{1/3}).                              \tag{4}
\]
Their product is integrable for every $n\geq1$. All lower mixed
distributions in (3) have an additional factor exponentially decaying
as $e^{-(n-k)L}$, so every energy summand is absolutely finite.
This is also the genuine decreasing-truncation energy: for
$w_j=\max(w,-j)$ the interface lies at $L_j\asymp j^3$.
Its top radial boundary mass is at most $Q'(-L_j)^n$, and its
energy contribution is $O(jL_j^{-2n/3})=o(1)$.
Lower mixed terms have the extra exponential factors. On the inner
truncated region the reference mass is exponentially small. Hence
all the truncated energy integrals converge to those just computed,
and $w\in\mathcal E^1$ with finite $E(w)<0$.

The FS reference distribution is $d(Q_0'^n)$. On the inner region
the relative probability density $f$ therefore satisfies
\[
 f=\frac{Q'^{n-1}Q''}{Q_0'^{n-1}Q_0''},\qquad
 \log f=nL-\left(\frac{2n}{3}+1\right)\log L+O(1).
                                                               \tag{5}
\]
The entropy tail is bounded by a constant times
$\int_1^\infty L^{-2n/3}\,dL$, finite for every $n\geq2$.
The remaining regions have finite entropy; the affine middle region
has zero top density. The universal negative-entropy bound controls
the negative part. Thus $H(w)<\infty$. This dimension threshold
explains why the same radial isolated-pole example on a curve has
infinite entropy.

\section{An actual nef model and its usual datum}
Let $H=\{Z_0=0\}$, which does not contain $p$. Blow up the smooth
codimension-two center $H\times\{0\}$ in $X\times\PP^1$:
\[
 \pi:\mathcal X\longrightarrow X\times\PP^1,
 \qquad \mathcal L=\pi_X^*\mathcal O(1)-E.
\]
The ordinary ideal is locally $(Z_0,\tau)$. The sections
$Z_0,\tau Z_1,\ldots,\tau Z_n$ have exactly this base ideal in the
appropriate bundle frames over the affine time chart, with no other
base locus there. On the blowup over that chart their divided sections
generate $\mathcal L$ and define a morphism
$g:\mathcal X|_{\mathbf A^1}\to\PP^n$, with
$\mathcal L|_{\mathbf A^1}=g^*\mathcal O(1)$.
Away from time zero the compactified model is the trivial product
with its original pulled-back bundle. Thus the actual rational bundle
is relatively semiample and relatively nef on the whole base.
It is not asserted to be relatively ample.

The central components are the strict product fiber $F_0$ and $E$,
both of time multiplicity one. The restriction of $\mathcal L$ to
$F_0\simeq X$ is trivial, and its top intersection on $E$ is one.
Indeed $E=\PP(\mathcal O_H(-1)\oplus\mathcal O_H)$ in the quotient
conormal convention. For $h=c_1(\mathcal O_H(1))$ and
$\xi=c_1(\mathcal O_E(1))$, we have
\[
 \xi^2+h\xi=0,\quad \pi_*\xi=1,\quad
 \int_E(h+\xi)^n=\sum_{k=1}^n\binom nk(-1)^{k-1}=1.
\]
The divisorial point of $E$ is $\operatorname{ord}_H$, and the blowup
relative canonical coefficient is one. The discrepancy formula
$(\kappa+1-b)/b$ is therefore one. The genuine model probability
measure and entropy are
\[
 \MA^{\rm NA}(\phi)=\delta_{\operatorname{ord}_H},
                         \qquad H^{\rm NA}(\phi)=1.      \tag{6}
\]
On the entire original NA carrier the actual section metric is
\[
                    \phi(v)=-\min\{v(H),1\}\in[-1,0].   \tag{7}
\]
Here $v(H)$ is the local effective Cartier-divisor order, with value
$+\infty$ at supports contained in $H$. Off $H$ one generator is a
unit; on $H$ another projective coordinate is a unit. This proves
(7) for supported semivaluations as well as function-field valuations.

For completeness set $T=\pi_X^*c_1(\mathcal O(1))$. The projective
bundle relation above gives $T^nE=0$ and
$T^{n+1-r}E^r=-1$ for $2\leq r\leq n+1$.
Consequently
\[
 E^{\rm NA}(\phi)=\frac{(T-E)^{n+1}}{n+1}
 =-\frac1{n+1}\sum_{r=2}^{n+1}(-1)^r\binom{n+1}r
 =-\frac n{n+1}.                                      \tag{8}
\]
All component multiplicities, line bundles and fixed factors are actual.

\section{The zero-start ray and its rigid absolute endpoint}
The associated actual total potential is
\[
 a_t=\log\frac{|Z_0|^2+e^{-t}\sum_{i=1}^n|Z_i|^2}
                       {\sum_{i=0}^n|Z_i|^2}.          \tag{9}
\]
It is bounded on each finite time annulus. On the blowup over the
affine time chart its total current is $g^*\omega_{\rm FS}$, whose
$(n+1)$st power vanishes.
It therefore solves the actual homogeneous geodesic equation.
Its datum is (7), by the divided-section compatibility. Bounded
geodesic energy is affine. At zero its derivative is minus the FS
average of $\sum_{i=1}^n|Z_i|^2/\sum_{i=0}^n|Z_i|^2$, hence
$E(a_t)=-nt/(n+1)$. Equation (8) and $R$'s actual energy-defect
criterion identify it with $r(0,\phi)$.

The ray decreases to $w_H$, which has ordinary curvature $[H]$.
Its nonpluripolar product is zero in every positive degree, in
particular its top raw mass is zero. Nevertheless it is a nontrivial
absolute $P$-model potential. To verify this directly, let
$h_0\leq0$ be any $\omega$-psh rooftop competitor below $w_H+C$.
Locally subtracting $\log|Z_0|^2$ from its total weight leaves a
psh function off $H$, bounded above there, hence with a psh extension.
Thus $\omega_{h_0}-[H]$ is a positive closed current of zero
cohomology class. Its integral against $\omega^{n-1}$ is zero,
so positivity makes it zero. On the compact connected $X$ this gives
$h_0=w_H+C'$; $\sup h_0\leq0$ forces $C'\leq0$.
Since $w_H$ itself is a competitor, $P_\omega[w_H]=w_H$.
The same reasoning proves the useful rigidity statement
\[
 h\in\operatorname{PSH}(X,\omega),\quad h\leq w_H+C
                      \quad\Longrightarrow\quad h=w_H+C'.       \tag{10}
\]

An optional direct entropy calibration also gives
\[
 H(a_t)=t+(n+1)\int_0^1n x^{n-1}
                     \log\bigl(x+e^{-t}(1-x)\bigr)\,dx,
 \qquad -\frac{n+1}{n}\leq H(a_t)-t\leq0.              \tag{11}
\]
Indeed its radial density relative to FS is
$e^{-nt}((1+r^2)/(1+e^{-t}r^2))^{n+1}$; use the transported
moment $x=e^{-t}r^2/(1+e^{-t}r^2)$, whose radial probability is
$n x^{n-1}dx$. This checks the probability and squared-log time.

\section{The singular-start ray loses its zero level}
The stationary ray $w$ has actual usual NA datum zero: on every
Gauss divisorial component centered at its pole, (2) is $O(t^{1/3})$
and hence has zero logarithmic-square coefficient; away from the
pole it is locally bounded. The canonical psh extension of those
divisorial data is zero. This uses the usual total-potential datum,
not pointwise division of the value $w(p)=-\infty$ by time.
Its energy slope is zero, so $R$ identifies $r(w,0)_t=w$.
Constant negative datum $-1$ gives $r(w,-1)_t=w-t$.
Order in both arguments and (7) imply
\[
                w-t\leq b_t=r(w,\phi)_t\leq w,
                \qquad b_t\leq a_t.                   \tag{12}
\]
At any point where $w$ is finite, the finite convex time function
$b_t$ is bounded above by $w$ and therefore is nonincreasing.
At $p$ all its values are $-\infty$. Its decreasing limit is
either a genuine $\omega$-psh function or identically $-\infty$.
If the limit were a nontrivial $h$, (9) and (12) would give
$h\leq w_H$ and $h\leq w$. Rigidity (10) makes
$h=w_H+C'$, finite at $p$, contradicting $w(p)=-\infty$.
Hence $\chi_b(0)=\inf_t b_t\equiv-\infty$, proving (1).

The ordered distance and exact energy identity in $R$ give
\[
 d_1(a_t,b_t)=E(a_t)-E(b_t)=-E(w)\quad\hbox{for every }t.          \tag{13}
\]
The singular ray is maximal and its initial entropy is finite, so
$R$ gives $H(b_t)/t\to H^{\rm NA}(\phi)=1$.
Every finite-time entropy is finite as well. Indeed this limit
supplies finite-entropy times tending to infinity. The finite-energy
Mabuchi functional is geodesically convex and is entropy plus
finite contracted energies \cite{BDL}. Its initial value and values
at those times are finite. Convexity makes its value finite at every
intermediate time; the contracted energies are finite on $\mathcal E^1$,
so entropy is finite there. The same assertion for $a$ follows from
(11). No entropy continuity under a varying initial point is asserted.

\section{What this example decides}
This is a finite-initial-entropy obstruction to absolute endpoint
invariance, even with bounded datum, finite discrepancy entropy,
constant finite ray distance and exact entropy-slope agreement.
Both zero-level raw masses are zero; those numbers do not identify
their $P$ data. A singularity metric is not applied to the identically
$-\infty$ level. The target model is nef and semiample; none of the
claims requires pretending it is relatively ample.

The example does not prove that no model sequence can recover a
prescribed singular-start ray or its relative endpoint. A sequence
retaining the singular initial potential is not excluded. Thus the
full unchanged original P19.2 existence problem remains open, and
this unreviewed supporting candidate changes no completion count.
It strengthens the earlier curve diagnostic by replacing its
infinite initial entropy with a finite one in every dimension $n\geq2$.

\begin{thebibliography}{9}
\bibitem{R} Complete established P19.1 mathematical theorem,
\emph{Entropy slopes of finite-energy K\"ahler rays: calibration,
recovery, and arbitrary starting points}, 6 October 2026.
Used only in its full fixed smooth compact K\"ahler signature.
\bibitem{BDL} R. Berman, T. Darvas and C. H. Lu,
\emph{Convexity of the extended K-energy and the large time behavior
of the weak Calabi flow}, Geometry \& Topology 21 (2017), 2945--2988,
\href{https://msp.org/gt/2017/21-5/gt-v21-n5-p08-s.pdf}{Theorem 4.7}.
In its operator convention use potential $u/(2\pi)$; currents and
probability entropy agree. Only finite contracted energies and
geodesic convexity are used.
\end{thebibliography}
\end{document}
